\section{What operator positivity the draft actually uses}
\label{sec:positivity}

\paragraph{MA-POS-001 (source audit and scope).}
A full-text search of the public draft for positivity, positive
operators, self-adjointness, and the associated projection and
singular-value arguments, followed by reading each relevant context,
gives the following inventory. The entries distinguish positivity
on an ambient Hilbert space from positivity compatible with the
arithmetic convolution action on the full zero-jet quotient.
The latter is the specific property addressed by GRH\_NIGHT
MT-M20 and MT-M23. None of the operator-positivity uses identified
in this inventory asserts or requires that property.

\begin{longtable}{p{0.25\linewidth}p{0.38\linewidth}p{0.28\linewidth}}
\hline
Exact draft locator & Operator and use & Does it require faithful
full-jet positivity?\\
\hline
\endhead
Remark 4.3.1.6, p.~176 & The bounded nonnegative Borel argument
$\varphi(A)$ of a unitary operator $A$ in the regular representation;
orthogonal spectral projections justify positivity and self-adjointness.
& No. This is the ordinary unitary regular representation.\\[3pt]
Corollary 4.3.3.4, pp.~184--185 & The compact-isotype averaging
operator is Haar volume times the orthogonal isotypic projector.
Peter--Weyl orthogonality proves positivity.
& No. It concerns the compact factor, whose action is unitary.\\[3pt]
Theorem 9.2.1.2(3), proof p.~405 & Three nonnegative indicator
multipliers are positive idempotents; hence their ranges give the
orthogonal splitting used in the Paley--Wiener diagram.
& No. Multiplication by an indicator is positive in an ordinary
weighted $L^2$ space.\\[3pt]
Remark 9.3.3.2, pp.~427--428 & A smooth real cutoff
$0\le\rho\le1$ gives a positive self-adjoint multiplier $M_\rho$.
& No. The inner product is the explicitly weighted integral.\\[3pt]
Remark 9.3.3.4, p.~428 & The absolute indicator cutoffs are positive
orthogonal projectors.
& No. The separate false injectivity assertion audited in MT-M34/35
does not make their positivity false.\\[3pt]
Definition 10.4.1.1, p.~456; Theorem 10.4.2.2, p.~461 & The
geometric and spectral orthogonal projections are positive by
definition. Self-adjointness of $M_\rho$ is used to pass a strong
cutoff limit through a trace-class operator.
& No. Orthogonality does not make the nonunitary translation
representation a star representation.\\[3pt]
Theorem 12.1.1.2, proof pp.~499--500 & An arbitrary nuclear
operator is decomposed into a linear combination of four positive
nuclear operators. The positive case is treated by spectral
expansion, then the general case by linearity.
& No. Positivity is an auxiliary reduction, not a claim that the
arithmetic convolution operator is positive.\\[3pt]
Lemma C.3.4.4(5), p.~699 & $|P_iF|$ is positive and finite rank;
its eigenvalues bound the trace norm of the finite-rank error.
& No. This is the absolute value of an arbitrary operator.\\[3pt]
Lemma C.3.5.3, pp.~699--700; Theorem C.3.6.1, pp.~703--704 &
The positive compact operator $|A|=(A^*A)^{1/2}$ supplies
singular values, the min--max formulas, and the singular-value
expansion.
& No. $|A|$ is positive regardless of whether $A$ is normal.\\[3pt]
Corollary C.3.5.5, pp.~702--703; Proposition C.3.6.5,
pp.~705--707 & Self-adjoint cutoff approximants convert strong
convergence into convergence in a Schatten ideal when composed
with an operator already in that ideal.
& No. The trace-class premise remains a separate premise.\\[3pt]
Section C.3.7, p.~707 & A kernel criterion is stated specifically
for positive operators; the Volterra example is used to explain
why the same criterion fails for arbitrary operators.
& No. The stated restriction is not replaced by a claim of
positivity for every geometric kernel.\\
\hline
\end{longtable}

\emph{Verification and limit of the audit.}
For every entry the surrounding definition, argument, and use of
the operator were read, rather than inferred from a keyword alone.
The multiplier entries follow directly from
$\langle M_\rho f,f\rangle=\int\rho|f|^2w\ge0$.
The compact-isotype entry acts by the positive Haar volume on its
isotype and by zero on the mutually orthogonal other isotypes.
The four-positive-operators reduction on p.~499 is the usual
decomposition of the real and imaginary self-adjoint parts into
positive and negative parts; all four remain trace class. Hence
it does not transfer positivity to the original arithmetic operator.
The absolute-value entries concern $A^*A$ and its nonnegative
square root, not the eigenvalues of $A$ itself. These observations
verify the stated scope of each use. They are not an independent
audit of every analytic theorem surrounding those uses.

Several other occurrences must not be added to this list as if
they were operator positivity. Sections 1.4.6 and 1.7.2,
pp.~30 and 35, discuss Weil's positive \emph{functional} and say
that the proof does not use it. The discussion on pp.~528 and
533 refers to an unfinished positivity criterion, including
``Theorem ??'' and the empty Theorem 13.3.1.2 on p.~522.
It supplies no additional positive arithmetic operator. The
positive Radon functional on p.~506 and the positive measure and
upper density on p.~635 are measure-theoretic positivity.
Likewise, ``positive cutoff'' in Definition 9.3.3.3 refers to
one side of a norm interval; operator positivity follows instead
from the nonnegative multiplier. We found no separate assertion
elsewhere in the checked positivity contexts that the faithful
full zero-jet action is a positive star representation. The audit
does not turn the absence of such an assertion into a proof of
the missing geometric trace identity.

\subsection{The weighted adjoint is different from arithmetic inversion}

\paragraph{MA-POS-002 (a direct weighted-space obstruction).}
Let $H=L^2(\mathbb R,w(x)\,dx)$ with the weight in (MA-TR-1),
where $\alpha>0$. For $U_yf(x)=f(x-y)$,
\[
 U_y^*g(x)=\frac{w(x+y)}{w(x)}g(x+y).
 \tag{MA-POS-1}
\]
In particular the adjoint is not $U_{-y}$ for $y\ne0$.
Even a compactly supported smooth convolution square need not
give a positive operator in this weighted space. More precisely,
if $0\ne f\in C_c^\infty(\mathbb R)$,
$f^\sharp(x)=\overline{f(-x)}$, and $g=f*f^\sharp$, then
$C_g$ is bounded on $H$ but is not self-adjoint and therefore
is not positive.

\begin{proof}
Substitute $x=u+y$ in
$\int f(x-y)\overline{g(x)}w(x)\,dx$ to obtain
(MA-POS-1). Boundedness of translations follows from the weight
estimate (M26.10), and integration over the compact support of
$g$ proves boundedness of $C_g$.
Conjugate to unweighted $L^2$ by the unitary map
$Vf=\sqrt w f$. The resulting integral kernel is
\[
 K(x,z)=g(x-z)\sqrt{\frac{w(x)}{w(z)}}.
\]
Since $g(-u)=\overline{g(u)}$, the adjoint kernel is
$g(x-z)\sqrt{w(z)/w(x)}$. Now
$g(0)=\int|f|^2>0$, so continuity makes $g(y)\ne0$ for all
sufficiently small $y$. For a positive such $y$ and $z>0$,
the weight is strictly increasing and $w(z+y)\ne w(z)$.
Thus the two continuous kernels differ on a nonempty open
rectangle away from the cusp at zero. Testing there against
two compactly supported smooth functions proves that the
bounded operators differ from their adjoints. Every positive
bounded operator is self-adjoint, proving the assertion.
\end{proof}

This calculation contains no zeta zeros. It explains precisely
why positivity of the source's multipliers and orthogonal
projections cannot be promoted to arithmetic star compatibility.
It is not an additional claim that the author made that promotion.
The source itself distinguishes the two projection orders in
(10.4.1.10)--(10.4.1.11), pp.~457--458.

\paragraph{MA-POS-003 (verification of the source's test-square observation, with its finite-sum form).}
Masullo already states the single-square observation in
\S1.7.2, p.~35. It is verified here and extended to finite
sums; it is not claimed as a newly discovered single-square
obstruction.
For an idele class group $C_F$ with Haar measure and involution
$f^\sharp(g)=\overline{f(g^{-1})}$, the test space
$\mathcal D^0(C_F)$ of MT-M30 contains no nonzero convolution
square $f*f^\sharp$ with $f\in\mathcal D(C_F)$.
\emph{Proof.} The idele class group is abelian, hence unimodular,
and its identity $e$ has norm one. Directly,
\[
 (f*f^\sharp)(e)=\int_{C_F}|f(g)|^2\,d^\times g.
\]
If the convolution square belongs to $\mathcal D^0$, the left
side is zero, so $f=0$ almost everywhere, and then identically
by continuity.\hfill$\square$
The same proof applies to a finite sum of such squares: its
value at $e$ is a sum of nonnegative squared norms. Thus an
identity restricted to norm-one-vanishing tests is not by itself
a positivity test on nonzero squares. MT-M30 provides an
ideal-valued correction that recovers equality of the relevant
linear functionals; it expressly does not say that this
correction preserves positivity.

\paragraph{MA-POS-004 (individual spectral projectors do not exhaust an operator).}
There is a separate literal convergence error in the positive-case
argument on p.~500. The paragraph immediately before (12.1.1.7)
defines $P_m$ as the projector onto
$\ker(T-\sigma_m(T)\operatorname{Id})$, with distinct singular
values listed once each, and then asserts
$P_mTP_m\to T$ in $\mathcal S_p$, $p=1,2$.
The paragraph preceding (12.1.1.12) similarly asserts convergence
of these individual traces to $\operatorname{Tr}T$.
Both assertions are false; cumulative projectors give the repair.

\emph{Proof and explicit counterexample.}
On $L^2(\mathbb R)$ choose an orthonormal Hermite basis
$(e_m)_{m\ge1}$ and define $Te_m=2^{-m}e_m$.
This is a positive trace-class operator of trace one with
distinct eigenvalues. Its individual eigenspace projector gives
\[
 \begin{aligned}
 \|T-P_mTP_m\|_1&=1-2^{-m}\longrightarrow1,\\
 \|T-P_mTP_m\|_2^2&=\tfrac13-4^{-m}\longrightarrow\tfrac13,\\
 \operatorname{Tr}(P_mTP_m)&=2^{-m}\longrightarrow0.
 \end{aligned}
\]
Thus none of the asserted individual limits holds.
Instead put $E_N=\sum_{m\le N}P_m$. For an arbitrary positive
trace-class $T$, its spectral expansion gives
\[
 \|T-E_NTE_N\|_1
   =\sum_{m>N}\sigma_m(T)\operatorname{rank}P_m\longrightarrow0,
\]
and hence convergence also in $\mathcal S_2$ and of traces.
The kernels $k_m$ of (12.1.1.7) remain the individual summands;
their partial sums are the kernels for $E_NTE_N$.
This repairs precisely the two uses of exhaustion.
\hfill$\square$
The source's later displayed partial sums show the intended
repair, but do not make its individual-projector sentences true.
This defect does not contradict the valid decomposition into
four positive operators, nor refute the corrected general
trace--kernel theorem. It is an additional source check, and
uses no zero-location premise.

\subsection{Faithfulness, multiplicity, and the determinant programme}

The exact known result remains MT-M20: positivity compatible with
the full convolution star action, in a faithful Hilbert completion
of the full zero-jet quotient, is equivalent to the corresponding
GRH assertion \emph{and simplicity}. MT-M23 applies
Hu--Kaneko--Martin--Schildkraut, Theorem 1.1, to nonabelian Galois
Dedekind zeta functions and rules out that faithful positive
completion there. These are carried as checked predecessor
results, not claimed as new theorems of this audit.

The adjective ``faithful'' must specify what structure it retains.
A single operator can be positive on a vector space that also
carries a nonunitary Jordan action: the identity operator is a
trivial example. The obstruction concerns simultaneous star
compatibility with the full arithmetic action. In a local jet
block this action contains a nonzero nilpotent extension; a
positive invariant Hermitian structure would make its centered
generator skew-adjoint, and hence diagonalizable, destroying
that extension. Positivity of one unrelated operator does not
have that implication.

For the actual determinant comparison, MT-M29 already supplies
the correct distinction. On the associated graded, repeated
zero ordinates may give repeated eigenvalues of a positive
trace-class inverse-square operator, with all multiplicities
retained. That construction discards the Jordan extensions. On
the positive-frequency part its Fredholm determinant is
\[
 \det(1-wB)=\prod_{\gamma>0}
              (1-w/\gamma^2)^{m_\gamma},
\]
with the central zero removed as specified in MT-M29.
This is compatible with multiple zeros; it is not a faithful
positive realization of the original full jet module.

In the inherited Grothendieck programme the precise archive
cross-references are cote~1, scans~079--080 (the determinant
coefficients, Corollaries 4--5, and the crossed-out product/trace
argument) and scans~111--112 (the genus-one discussion), as
audited in ALPOGE\_BRIDGES BR-TR-01/02 and F-SM-11/12.
Those scans were not newly supplied or freshly inspected here.
Passing from a finite filtration to its graded object preserves
its class in $K_0$, its traces, and the corresponding determinant
coefficients; this is why a $K$-theoretic or $\lambda$-ring
identity cannot on its own certify faithfulness of the jet action.
The algebraic formalism is SGA~6, Expos\'e~V, \S\S2 and 7,
especially (7.1.1), as placed in GRH\_NIGHT AG-LAMBDA.
The analytic positivity of a Hilbert completion is additional
structure, not a consequence of that formalism or of an
unconstructed cohomological realization over
$\operatorname{Spec}\mathbb Z$.

MA-POS-001 therefore clears a specific misdiagnosis: the known
full-jet obstruction is not a blanket refutation of every
positive cutoff used by the draft. MA-POS-002 identifies the
actual distinction of adjoints, and MA-POS-003 identifies the
restriction on the test cone. Neither proves the absent
global-to-local trace transport.
